\section{Context Menu And Saved Settings}\label{sec:settings}

Right-click the module to find these options under \textbf{Render Settings}:
\begin{itemize}
  \item \textbf{AC-coupled} removes DC offset from newly captured inputs
  using a DC-blocking filter. It is on by default and can also affect very
  low frequencies. Turn it off when inspecting DC or slow signals, and set
  Slope to zero to avoid frequency weighting near DC. This setting does not
  reprocess history.
\end{itemize}

Palette and intensity-scale selection are on the color screen rather than
in the context menu. New instances and reset use Decibels; old patches and
presets with no valid intensity mode load in Linear to preserve their
colors. Missing endpoints load as $-90$/$0\,\mathrm{dB}$ and 0/100\%, even
when loading into an edited module.

Panel dropdowns and AC-coupling changes support Rack's Undo and Redo commands. Patches and module
presets save the panel settings, Run state, AC-coupling choice, color map,
and intensity mode.
They do not save the recorded spectrogram. A patch saved while frozen reopens
with no recorded history; enable Run to acquire a new image.

Rack's module initialization resets the controls to their defaults, enables
Run and AC coupling, selects Magma and Decibels, restores Floor/Ceil to
$-90$/$0\,\mathrm{dB}$ and Linear to 0/100\%, and clears the history. Changing Rack's
sample rate changes the available frequency range, bin spacing, and sweep
duration. Existing history is not cleared automatically and is displayed
using the new frequency mapping; allow a full sweep to replace it before
comparing frequencies at the new sample rate.
